\documentclass[12pt]{article}

% Pachete necesare pentru codificarea caracterelor și diacritice
\usepackage[utf8]{inputenc}
\usepackage[romanian]{babel}

% Pachete matematice avansate
\usepackage{amsmath}
\usepackage{amsfonts}
\usepackage{amssymb}
\usepackage{amsthm}

% Setarea marginilor paginii
\usepackage[margin=1in]{geometry}

\title{Seminar 2: Criptografie -- Rezolvare Teme}
\author{}
\date{} % Această linie goală elimină complet data și spațiul ei sub titlu

\begin{document}

\maketitle


\subsection*{Exercitiul 2}
\textit{Găsiți numărul maxim de pași pentru algoritmul lui Euclid. Explicați în ce context se obține.} 


\subsection{Schiță de Demonstrație}
Fie resturile succesive generate de algoritmul lui Euclid aplicat numerelor $a > b$:
\begin{align*}
a &= b \cdot q_1 + r_1, \quad 0 < r_1 < b \\
b &= r_1 \cdot q_2 + r_2, \quad 0 < r_2 < r_1 \\
  &\ \ \vdots \\
r_{k-2} &= r_{k-1} \cdot q_k + r_k, \quad 0 < r_k < r_{k-1} \\
r_{k-1} &= r_k \cdot q_{k+1} + 0
\end{align*}
Algoritmul s-a oprit după exact $k+1$ pași, iar ultimul rest nenul este $r_k \ge 1$. Deoarece resturile sunt strict descrescătoare și toate câturile $q_i \ge 1$, deducem inegalitățile:
\begin{itemize}
    \item $r_k \ge 1 = F_2$
    \item $r_{k-1} \ge 2 \cdot r_k \ge 2 = F_3$ (deoarece $r_{k-1} > r_k \implies q_{k+1} \ge 2$)
    \item $r_{k-2} = r_{k-1} \cdot q_k + r_k \ge r_{k-1} + r_k \ge F_3 + F_2 = F_4$
\end{itemize}

Prin inducție matematică, se demonstrează ușor că pentru orice pas $i$ avem:
\[
r_{k-i} \ge F_{i+2}
\]
Pentru primul pas al algoritmului (unde $i = k$), relația devine:
\[
b = r_0 \ge F_{k+2}
\]
Folosind proprietatea de creștere exponențială a șirului lui Fibonacci bazată pe Secțiunea de Aur ($\phi = \frac{1+\sqrt{5}}{2} \approx 1.618$), se știe că $F_{k+2} > \phi^k$. Prin urmare, avem:
\[
b > \phi^k
\]
Aplicând logaritmul în baza 10 în ambele părți ale inegalității obținem:
\[
\log_{10}(b) > k \cdot \log_{10}(\phi)
\]
Deoarece $\log_{10}(\phi) \approx \log_{10}(1.618) > 0.2$, rezultă că:
\[
\log_{10}(b) > \frac{k}{5} \implies k < 5 \cdot \log_{10}(b)
\]
Cum numărul total de pași este $k+1$, iar inegalitatea precedentă este strictă pentru mărgini întregi, se confirmă limita superioară impusă de Teorema lui Lamé.

\end{document}